\subsection{A projective family with the prescribed singularities}
\label{sec:compactification}

The following construction allows the local invariant to be compared by
ordinary intersection numbers on projective surfaces.  It preserves each
individual analytic right type; no simultaneous analytic coordinate
change is needed.

\begin{lemma}[Compactification with one singular point]
\label{lem:compactification}
Let $C$ be a smooth connected complex algebraic curve with marked point
$0$.  Let $g_c(x_1,x_2,x_3)$ be polynomials of bounded degree $l$, with
coefficients regular on $C$, satisfying $g_c(0)=0$ and
$d_xg_c(0)=0$.  Suppose the critical point at the origin is isolated and
has the same finite Milnor number $\mu$ for every $c\in C$.
After replacing $C$ by a Zariski neighborhood of $0$, there is a flat
projective hypersurface family $X\subset\mathbb P^3\times C$ of some
fixed degree $d$, with section $\Gamma$, such that each fibre is integral
and normal, smooth outside $\Gamma(c)$, and its germ at $\Gamma(c)$
admits a defining function analytically right-equivalent to $g_c$.
Furthermore,
\begin{align}
 K_{X_c}^2&=d(d-4)^2,\label{eq:compact-canonical-square}\\
 \chi_c(X_c\setminus\{\Gamma(c)\})
   &=d(d^2-4d+6)-\mu-1.\label{eq:compact-euler}
\end{align}
The total space $X$ is normal and Gorenstein.  These conclusions remain
valid after a finite change of smooth base curves.
\end{lemma}

\begin{proof}
Choose $d>l$ with $d\geq\mu+2$ and homogeneous coordinates
$[x_1:x_2:x_3:w]$ on $\mathbb P^3$.  Homogenize $g_c$ to degree $d$,
obtaining $G_c(x,w)$, and set
\[
 X_c=\{G_c(x,w)+Q_d(x)=0\},
 \qquad \Gamma(c)=[0:0:0:1],
\]
where $Q_d$ is a general homogeneous form of degree $d$ in the three
spatial coordinates, independent of $c$.

Since $Q_d\in\mathfrak m^{\mu+2}$, Lemma~\ref{lem:individual-jets}
shows, for each $c$ separately, that adding $Q_d$ preserves the analytic
right type of $g_c$.  This finite-determinacy and high-degree-tail
compactification is also the construction used in
\cite[Lemma~2.5]{Steenbrink1985}.

At $c=0$, the linear system spanned by $G_0$ and all forms $Q_d(x)$
has base locus exactly $\{[0:0:0:1]\}$.  The characteristic-zero
Bertini theorem \cite[Theorem~4.1]{KleimanBertini} therefore makes a
general member smooth elsewhere;
its nonzero coefficient of $G_0$ can be normalized to one.  One can
also see the required Bertini assertion directly.  On
$U=\mathbb P^3\setminus\{[0:0:0:1]\}$, the incidence variety
\[
 \{(z,Q):G_0(z)+Q(z)=0\}\subset U\times H^0(\mathbb P^2,\mathcal O(d))
\]
is an affine bundle over $U$, since evaluation of the forms is
surjective at every $z\in U$.  It is smooth.  Its generic fibre over
the coefficient space is regular, being obtained by localization, and
hence smooth over its characteristic-zero function field.  The
nonsmooth locus has constructible image missing the generic point,
so a nonempty open set of choices has smooth fibres.

There are no other singularities for complex parameters sufficiently
near $0$.  Indeed, otherwise projectivity gives a sequence of singular
points converging to the unique singular point of $X_0$.  They are then
in the chart $w=1$.  Conservation of the spatial-gradient number in a
small ball contradicts their existence, since the origin already
accounts for the full constant number $\mu$.  The parameters with an
additional singular point form a constructible subset of $C$.  A
constructible subset of an algebraic curve which misses an analytic
neighborhood of $0$ is finite; removing it gives the required Zariski
neighborhood.

Every resulting projective hypersurface has only one singular point.
It is reduced and irreducible: a repeated component or the intersection
of two distinct surface components would give a positive-dimensional
singular locus.  The hypersurface is Cohen--Macaulay and regular in
codimension one, hence normal.  Since the defining equation is nonzero
on every ambient fibre, it is a relative effective Cartier divisor and
the family is flat.  Outside the section it is smooth over $C$.
The total hypersurface is Gorenstein and regular in codimension one,
so it too is normal.

Adjunction gives $\omega_{X_c}=\mathcal O_{X_c}(d-4)$ and therefore
\eqref{eq:compact-canonical-square}.  A small global smoothing changes
only a Milnor neighborhood of the singular point: outside it the
family is a proper submersion with transverse boundary.  Replacing
the contractible singular neighborhood by its Milnor fibre increases
Euler characteristic by $\mu$.  A smooth degree-$d$ surface in
$\mathbb P^3$ has Euler characteristic $d(d^2-4d+6)$, as follows from
its tangent-normal exact sequence and the Gauss--Bonnet formula.
This proves \eqref{eq:compact-euler}.  Such a smoothing can, for example,
be obtained by subtracting $\varepsilon w^d$; near the isolated
critical point this is the usual Milnor smoothing, and compactness
preserves smoothness away from that neighborhood.

A change of base leaves the individual projective surfaces and their
germs unchanged.  Over a smooth new curve the relative hypersurface
description, flatness, smoothness off the section, and the same
codimension-one argument give all the remaining assertions again.
\end{proof}
